\title{Controlling Text-to-Image Diffusion by Orthogonal Finetuning}

\begin{abstract}
Recent advances in large text-to-image diffusion models have enabled the synthesis of highly realistic images from textual descriptions. However, efficiently steering these models for various downstream applications remains a significant challenge. In response, we present a theoretically grounded finetuning approach, Orthogonal Finetuning~(OFT), designed to tailor text-to-image diffusion models for specific tasks. Distinct from existing approaches, OFT is theoretically shown to maintain hyperspherical energy, which encapsulates pairwise relationships among neurons on the unit hypersphere. Preserving this quantity is shown to be essential for maintaining the semantic fidelity of generation in text-to-image diffusion models. To further enhance the robustness of the finetuning process, we introduce Constrained Orthogonal Finetuning (COFT), incorporating a radius constraint on the hypersphere. Our study concentrates on two prominent finetuning scenarios for text-to-image models: subject-driven generation, where the aim is to synthesize images of a specific entity using a limited set of its images and a text description; and controllable generation, wherein the model is provided with auxiliary control cues in addition to textual input. Experimental evaluation demonstrates that the OFT framework offers notable improvements over prior techniques in both image quality and speed of convergence.
\end{abstract}

\section{Introduction}

Modern text-to-image diffusion models~\cite{saharia2022photorealistic,ramesh2022hierarchical,rombach2022high} have set new benchmarks in producing visually compelling images guided by textual input. Nonetheless, using text alone for control can lead to ambiguity, often falling short in granting precise, fine-level manipulation of generated content. To tackle this issue, our work investigates two central text-to-image generation paradigms:

\begin{itemize}[leftmargin=*,nosep]

    \item \textbf{Subject-driven generation}~\cite{ruiz2023dreambooth}: With only a small set of reference images for a particular subject, the challenge is to synthesize new visuals depicting the same subject in novel scenarios, informed by a text prompt.
    \item \textbf{Controllable generation}~\cite{zhang2023adding,mou2023t2i}: By introducing auxiliary signals (such as canny edges or segmentation masks), this task seeks to produce images that align with both the given control input and accompanying textual description.
\end{itemize}

Both tasks fundamentally address the challenge of finetuning text-to-image diffusion models while retaining the generative capabilities acquired during pretraining. We outline the main requirements for a robust finetuning approach as follows: (1) \emph{training efficiency}, characterized by a reduced number of trainable parameters and fewer training epochs, and (2) \emph{generalizability preservation}, which entails maintaining the quality and diversity of the model's generative outputs. Common finetuning practices include either incrementally adjusting neuron weights using a small learning rate (\emph{e.g.}, \cite{ruiz2023dreambooth}) or introducing a compact module with separately parameterized neuron weights (\emph{e.g.}, \cite{hulora2022,zhang2023adding}). Although these techniques are straightforward, they do not offer guarantees for retaining the generative strengths established during pretraining. Furthermore, the development of effective finetuning procedures and the selection of key hyperparameters, such as training duration, often lack a systematic basis. A principal challenge in this area is the absence of metrics capable of rigorously measuring how well the finetuned model preserves the generative potential inherited from pretraining. Present finetuning strategies typically presume that a minimized Euclidean distance between the weights of the pretrained and finetuned models reflects superior preservation of the pretrained characteristics. As a result, practitioners frequently employ exceedingly low learning rates. However, while this hypothesis might sometimes prove valid, we contend that Euclidean proximity alone fails to adequately describe the extent of semantic consistency between the two models. Consequently, incorporating a more structural criterion to assess the discrepancy between the pretrained and finetuned models could substantially enhance both the safeguarding of pretraining abilities and the overall stability of the finetuning process.

Building upon the empirical finding that hyperspherical similarity effectively captures semantic relationships~\cite{liu2017deep,liu2018decoupled,chen2020angular}, our approach leverages hyperspherical energy~\cite{liu2018learning} to describe the pairwise connections among neurons. Hyperspherical energy, formulated as the aggregate of hyperspherical similarities (such as cosine similarity) between all neuron pairs in a given layer, serves as a measure of how uniformly neurons are distributed over the unit hypersphere~\cite{liu2021learning}. We propose that optimal finetuned models should exhibit negligible deviation in hyperspherical energy relative to their pretrained counterparts. While one straightforward method might be to regularize the finetuning process so the hyperspherical energy remains constant, this approach does not ensure the difference in hyperspherical energy is effectively minimized. Consequently, we utilize a key invariance of hyperspherical energy—specifically, that the pairwise hyperspherical similarities remain invariant when a common orthogonal transformation is applied to all neurons. Drawing on this invariance, we introduce Orthogonal Finetuning (OFT), a technique to adapt large text-to-image diffusion models for downstream tasks while maintaining their hyperspherical energy. This method involves learning a shared orthogonal transformation per layer, thereby preserving the angles between neurons. Interpreted differently, OFT can be regarded as redefining the canonical basis for the neurons within each layer. Moreover, considering that maintaining a smaller Euclidean distance between the finetuned and pretrained models aids in preserving pretraining effectiveness, we further introduce a variant termed Constrained Orthogonal Finetuning (COFT), which restricts the finetuned parameters to lie on a hypersphere of fixed radius centered at the pretrained neurons.

The rationale behind employing orthogonal transformations for finetuning neurons is partly inspired by insights from the 2D Fourier transform, a technique that decomposes an image into magnitude and phase spectra. Notably, the phase spectrum—representing angular relationships between the input and the basis—encapsulates most of the image's semantic content. For instance, if one combines an image's phase spectrum with a random magnitude spectrum, it is still possible to reconstruct a version of the original image that maintains its semantic meaning. This observation underscores that altering neuron directions is fundamentally responsible for semantic changes in generated images, an objective central to both subject-driven and controllable image synthesis. Nevertheless, granting too much freedom in adjusting neuron directions can compromise the pretrained model’s generative abilities. To manage this, we advocate preserving the pairwise angles between neurons, motivated by the notion that these angular relationships encode essential knowledge within neural networks. This leads to our proposal of applying a layer-wise orthogonal transformation—shared across all neurons in a layer—to keep the hyperspherical energy invariant.

We also derive conceptual motivation from orthogonal over-parameterized training~\cite{liu2021orthogonal}, which initiates classification neural networks by subjecting a randomly initialized architecture to orthogonal transformations during training. Since a randomly initialized network is expected to have low hyperspherical energy, \cite{liu2021orthogonal} endeavors to maintain this property throughout training, as reduced hyperspherical energy is associated with improved generalization in classification tasks~\cite{liu2018learning,lin2020regularizing}. Their findings indicate that orthogonal transformations suffice to build well-generalizing classifiers from scratch. By contrast, our emphasis is on finetuning pretrained text-to-image diffusion models, aiming for enhanced control and superior generative outcomes. We highlight two principal distinctions between our OFT method and the approach in \cite{liu2021orthogonal}. First, unlike \cite{liu2021orthogonal}, which targets the minimization of hyperspherical energy, our OFT approach seeks to retain the pretrained model's original hyperspherical energy, thereby safeguarding its inherent semantic organization. In finetuning diffusion models, driving the hyperspherical energy lower may undermine foundational semantic integrity. Second, OFT is formulated to minimize divergence from the pretrained network—a constraint absent from the method in \cite{liu2021orthogonal}. Effective finetuning therefore requires balancing adaptation and model stability, and we argue that OFT provides an optimal trade-off in this context. The main contributions of our work are summarized below:

\begin{itemize}[leftmargin=*,nosep]

    \item We introduce Orthogonal Finetuning (OFT), a new approach for finetuning text-to-image diffusion models that enhances controllability. To address stability concerns, we further design a constrained version of OFT that bounds the angular divergence from the original pretrained weights.
    \item In contrast to prior finetuning techniques, OFT updates model parameters in a manner that mathematically guarantees conservation of hyperspherical energy—a property we empirically show to correlate strongly with the preservation of generative semantics from the pretrained model.
    \item Our method is evaluated on two primary tasks: subject-driven generation and controllable generation. Through extensive empirical experiments, we show that OFT yields marked gains over existing methods in terms of sample quality, finetuning stability, and convergence rate. Additionally, OFT is more sample-efficient, achieving effective convergence using far fewer training images and epochs.
    \item For the controllable generation setting, we propose a novel control consistency metric that quantifies controllability by comparing the original control input and a signal inferred from the model’s generated output. Our experiments further demonstrate the superior controllability achieved by OFT, as validated by this new metric.
\end{itemize}

\section{Related Work}

\textbf{Text-to-image diffusion models}. Significant advancements~\cite{saharia2022photorealistic,ramesh2022hierarchical,rombach2022high,gu2022vector,nichol2022glide} have been realized in text-to-image synthesis, mainly propelled by innovations in diffusion-based generative models~\cite{ho2020denoising,song2021score,song2021denoising,dhariwal2021diffusion} as well as progress in joint vision-language representation learning~\cite{su2020vlbert,lu2019vilbert,radford2021learning,wang2021simvlm,li2022blip,li2023blip,alayrac2022flamingo,schuhmann2022laion}. Both GLIDE~\cite{nichol2022glide} and Imagen~\cite{saharia2022photorealistic} implement diffusion modeling directly in pixel space. While GLIDE leverages joint training of the text encoder and diffusion model using paired text-image examples, Imagen instead employs a static, pretrained text encoder throughout training. Meanwhile, Stable Diffusion~\cite{rombach2022high} and DALL-E2~\cite{ramesh2022hierarchical} train their diffusion models in a latent space. The latent space in Stable Diffusion is defined via a VQ-GAN~\cite{esser2021taming} visual codebook, whereas DALL-E2 utilizes CLIP~\cite{radford2021learning} embeddings to create a shared representation for images and text. In addition to diffusion models, other paradigms such as generative adversarial networks~\cite{reed2016generative,zhang2017stackgan,xu2018attngan,li2019controllable} and autoregressive approaches~\cite{ramesh2021zero,ding2021cogview,wu2022nuwa,yuscaling} have also been explored within the domain of text-driven image generation. OFT is designed as a model-agnostic finetuning strategy and is compatible with any text-to-image diffusion architecture.

\textbf{Subject-driven generation}. Strategies to restrict changes to the subject, as in \cite{avrahami2022blended,nichol2022glide}, involve providing a user-specified mask as additional conditioning information. Inversion approaches~\cite{choi2021ilvr,dhariwal2021diffusion,ramesh2022hierarchical,gal2022image} offer mechanisms to edit contextual aspects of an image while maintaining the integrity of the central subject. The technique in \cite{hertz2022prompt} enables both local and global modifications without reliance on input masks. However, these previous solutions struggle to accurately retain subject-specific, identity-related features. The Pivotal Tuning framework~\cite{roich2022pivotal} addresses this by finetuning the generator around an inverted latent code, employing regularization to preserve identity. In a similar vein, \cite{nitzan2022mystyle} derives a personalized generative prior based on a collection of individual face images. While \cite{casanova2021instance} permits the generation of varied instances of an object, this comes at the risk of eroding instance-unique traits. DreamBooth~\cite{ruiz2023dreambooth}, through the use of a custom token and a handful of subject-specific images, optimizes a reconstruction loss and employs a class-prior preservation loss during fin

\textbf{Controllable generation}. Image-to-image translation represents a subclass of controllable generation, where much of the prior research relies on conditional generative adversarial networks~\cite{isola2017image,zhu2017toward,wang2018high,park2019semantic,choi2018stargan}. Additionally, diffusion models have been increasingly utilized in this context~\cite{saharia2022palette,voynov2022sketch,wang2022pretraining}. Recently, ControlNet~\cite{zhang2023adding} demonstrated an approach to regulating a pretrained diffusion model through finetuning with supplementary control inputs, achieving notable advances in controllable generation. Similarly, the concurrent method T2I-Adapter~\cite{mou2023t2i} refines a pretrained diffusion model to further strengthen generation controllability. In line with the experimental framework of \cite{zhang2023adding,mou2023t2i}, we employ OFT for finetuning pretrained diffusion models, which delivers improved controllability using reduced amounts of training data and requiring fewer parameters for finetuning. Importantly, OFT does not add extra computational expense during inference.

\textbf{Model finetuning}. The practice of finetuning large pretrained models for downstream applications has seen substantial adoption in recent years~\cite{devlin2018bert,brown2020language,he2022masked}. Within this broader trend, parameter-efficient adaptation strategies (\emph{e.g.}, \cite{hulora2022,houlsby2019parameter,pfeiffer2020adapterfusion}) have received significant attention in natural language processing. The work most related to OFT is LoRA~\cite{hulora2022}, which leverages the assumption of a low-rank structure in the incremental weight updates performed during finetuning. In contrast, OFT employs a layer-wise shared orthogonal transformation that modifies neuron weights multiplicatively and, importantly, mathematically guarantees the preservation of angles between neuron vectors within a layer, contributing to increased stability.

\section{Orthogonal Finetuning}

\subsection{Why Does Orthogonal Transformation Make Sense?}

To motivate the adoption of orthogonal transformation for finetuning in text-to-image diffusion models, we break down the rationale into two components: (1) the motivation for adjusting the angular orientation (\emph{i.e.}, the direction) of neuron vectors, and (2) the justification for using orthogonal transformations as the mechanism for finetuning these angles.

Addressing the first question, we take cues from the empirical findings of \cite{liu2018decoupled,chen2020angular}, which suggest that differences in angular features serve as a robust indicator of semantic disparity. SphereNet~\cite{liu2017deep} provides evidence that constraining all neurons to lie on a unit hypersphere does not diminish a neural network’s representational capacity, and may, in fact, enhance its generalization, suggesting that neuron direction is key for encoding crucial data features. To more explicitly illustrate the significance of neuron angles, we design a simple illustrative experiment shown in Figure~\ref{fig:toy_ae}. Here, a basic convolutional autoencoder is trained using flower image data, with the training phase employing the standard inner product to compute the feature map (where $z$ represents the output for each kernel convolution between $\bm{w}$ and the windowed input $\bm{x}$). In the evaluation phase, we compare three methods for producing the feature map: (a) inner product (mirroring the training setup), (b) the feature magnitude, and (c) the angle information. Results depicted in Figure~\ref{fig:toy_ae} reveal that using angular information allows almost perfect reconstruction of the original images, whereas relying on neuron magnitudes yields no meaningful information. It is important to note that cosine activation (c) is not used during training, which strictly utilizes inner product. These findings highlight that neuron directionality is primarily responsible for retaining the semantic content of input images, indicating that adjusting neuron directions is likely to be the most effective means for altering semantic information through finetuning.

In response to the second question, the most straightforward technique for finetuning neuron direction is to apply a joint rotation or reflection to all neurons within the same layer, which introduces an orthogonal transformation. Although it is possible to consider more adaptive angular transformations, where neurons are rotated by distinct angles, our results show that orthogonal transformations strike an optimal balance between adaptability and regularization. Supporting this perspective, \cite{liu2021orthogonal} demonstrates that orthogonal transformations possess ample expressive power in the context of neural network training. To empirically justify our approach, we conduct a regularization analysis through a controllable image generation task, following the ControlNet~\cite{zhang2023adding} configuration. Results in Figure~\ref{fig:ortho} show that our standard OFT yields consistent and accurate control by epoch 20. In contrast, the absence of an orthogonality constraint in OFT leads to a total failure in producing convincing images or exhibiting control effects. This experiment confirms the critical role of orthogonality in constraining OFT’s transformations.

\subsection{General Framework}

In standard finetuning procedures, gradient descent with a reduced learning rate is employed to adjust the model parameters—sometimes targeting only specific layers. This conservative update encourages model weights to stay close to those of the pretrained network, effectively corresponding to an implicit minimization of $\thickmuskip=2mu \medmuskip=2mu \| \bm{M} - \bm{M}^0\|$, where $\bm{M}$ and $\bm{M}^0$ represent the finetuned and initial pretrained weights, respectively. While such a constraint is intuitively reasonable, it may lack the necessary rigidity when finetuning large-scale models. LoRA responds to this limitation by imposing an additional constraint of low rank on the parameter update: specifically, $\thickmuskip=2mu \medmuskip=2mu \text{rank}(\bm{M} - \bm{M}^0)=r'$ for a small $r'$. In contrast, OFT targets the similarity of pairwise neurons via the constraint $\thickmuskip=2mu \medmuskip=2mu \| \text{HE}(\bm{M})-\text{HE}(\bm{M}^0) \|=0$, where $\text{HE}(\cdot)$ is the hyperspherical energy for a weight matrix.

To illustrate, take the case of a fully connected layer described by $\thickmuskip=2mu \medmuskip=2mu \bm{W}=\{\bm{w}_1,\cdots,\bm{w}_n\}\in\mathbb{R}^{d\times n}$, with each column vector $\bm{w}_i\in\mathbb{R}^d$ signifying a neuron ($\bm{W}^0$ denoting its pretrained counterpart). The output $\thickmuskip=2mu \medmuskip=2mu \bm{z}\in\mathbb{R}^n$ is determined by $\thickmuskip=2mu \medmuskip=2mu \bm{z}=\bm{W}^\top\bm{x}$ for a given input $\thickmuskip=2mu \medmuskip=2mu \bm{x}\in\mathbb{R}^d$. Under OFT, the focus shifts to reducing the difference in hyperspherical energies between the trained and original models:

\begin{equation}\label{eq:oft_principle}
\footnotesize
\min_{\bm{W}} \left\| \text{HE}(\bm{W})-\text{HE}(\bm{W}^0)\right\|~~\Leftrightarrow~~\min_{\bm{W}} \bigg{\|} \sum_{i\neq j} \|\hat{\bm{w}}_i-\hat{\bm{w}}_j\|^{-1}-\sum_{i\neq j} \|\hat{\bm{w}}^0_i-\hat{\bm{w}}^0_j\|^{-1}\bigg{\|}
\end{equation}

Here, the normalized neuron is $\thickmuskip=2mu \medmuskip=2mu \hat{\bm{w}}_i:=\bm{w}_i/\|\bm{w}_i\|$, and the hyperspherical energy for a layer $\bm{W}$ is given by $\thickmuskip=2mu \medmuskip=2mu \text{HE}(\bm{W}):= \sum_{i\neq j} \|\hat{\bm{w}}_i-\hat{\bm{w}}_j\|^{-1}$. It is straightforward to verify that the minimum achievable value for Eq.~\eqref{eq:oft_principle} is zero. This minimum is realized when $\bm{W}$ is related to $\bm{W}^0$ solely through a rotation or reflection, expressed as $\thickmuskip=2mu \medmuskip=2mu \bm{W}=\bm{R}\bm{W}^0$, with the orthogonal transformation $\thickmuskip=2mu \medmuskip=2mu \bm{R}\in\mathbb{R}^{d\times d}$ (where $\det(\bm{R})=1$ or $-1$ signifies rotation or reflection, respectively). This encapsulates the primary motivation behind OFT, which proposes finetuning a model by optimizing orthogonal matrices shared across layers to transform neuron weights. Concretely, OFT learns an orthogonal matrix $\thickmuskip

\begin{equation}\label{eq:oft_general}
    \footnotesize
    \bm{z} = \bm{W}^\top\bm{x} = (\bm{R} \cdot \bm{W}^0)^\top\bm{x},~~~\text{s.t.}~\bm{R}^\top\bm{R} = \bm{R}\bm{R}^\top = \bm{I}
\end{equation}

Here, $\bm{W}$ is the weight matrix after OFT-finetuning, and $\bm{I}$ represents the identity matrix. The structure of OFT is depicted in Figure~\ref{fig:block}. As with the zero initialization approach used in LoRA, it is essential for OFT to begin fine-tuning from the initial pretrained parameters $\bm{W}^0$. To accomplish this, the orthogonal matrix $\bm{R}$ is initialized to the identity matrix, ensuring the model's parameters are initially unchanged from their pretrained state. Maintaining orthogonality of $\bm{R}$ throughout training can be achieved via differentiable orthogonalization methods introduced in \cite{liu2021orthogonal,lezcano2019cheap}. We later elaborate on methods for maintaining this constraint efficiently in practice.

\subsection{Efficient Orthogonal Parameterization}\label{sect:eff_ortho}

Traditional approaches for constructing orthogonal matrices, such as the Gram-Schmidt process, are differentiable but often prove computationally prohibitive for large matrices~\cite{liu2021orthogonal}. To improve efficiency, we employ the Cayley parameterization, which generates orthogonal matrices in a more tractable manner. In this scheme, the orthogonal matrix is given by $\thickmuskip=2mu \medmuskip=2mu \bm{R} = (\bm{I} + \bm{Q})(I - \bm{Q})^{-1}$, where $\bm{Q}$ denotes a skew-symmetric matrix with the property $\thickmuskip=2mu \medmuskip=2mu \bm{Q} = -\bm{Q}^\top$. This method introduces a minor constraint: only orthogonal matrices with determinant $1$—those within the special orthogonal group—can be constructed, but empirical evidence suggests this restriction does not impair effectiveness. Nevertheless, representing a full orthogonal matrix $\bm{R}$ becomes increasingly inefficient as the dimension $d$ grows. To mitigate this, we propose a block-diagonal structure for $\bm{R}$ composed of $r$ separate blocks, each block being an independent orthogonal matrix, resulting in the following structure:

\begin{equation}
    \footnotesize
    \bm{R} = \text{diag}(\bm{R}_1, \bm{R}_2, \cdots, \bm{R}_r) = \begin{bmatrix}
    \bm{R}_1 \in \text{O}(\frac{d}{r}) &  &\\
    &\ddots & \\
    & & \bm{R}_r \in \text{O}(\frac{d}{r})
    \end{bmatrix} \in \text{O}(d)
\end{equation}

where $\text{O}(d)$ stands for the orthogonal group of $d$ dimensions, with $\thickmuskip=2mu \medmuskip=2mu \bm{R}\in\mathbb{R}^{d\times d}$ and $\thickmuskip=2mu \medmuskip=2mu \bm{R}_i\in\mathbb{R}^{d/r\times d/r},\forall i$ representing orthogonal matrices. When $\thickmuskip=2mu \medmuskip=2mu r=1$, the block-diagonal structure collapses to a regular, fully unconstrained orthogonal matrix. For a general $d\times d$ orthogonal matrix, there are $\thickmuskip=2mu \medmuskip=2mu d(d-1)/2$ independent parameters, corresponding to $\mathcal{O}(d^2)$ complexity. For an $r$-block diagonal orthogonal matrix, this reduces to $\thickmuskip=2mu \medmuskip=2mu d(d/r-1)/2$ parameters and thus a complexity of $\thickmuskip=2mu \medmuskip=2mu \mathcal{O}(d^2/r)$. Furthermore, one may decrease the parameter count by sharing block matrices, i.e., $\thickmuskip=2mu \medmuskip=2mu\bm{R}_i=\bm{R}_j,\forall i\neq j$, which yields a parameter complexity of $\mathcal{O}(d^2/r^2)$. Despite these methods of increasing parameter efficiency, the matrix $\bm{R}$ continues to satisfy the orthogonality property, ensuring the preservation of hyperspherical energy.

We next compare OFT and LoRA in terms of parameter efficiency. For LoRA with a rank $r'$, the number of tunable parameters is $\thickmuskip=2mu \medmuskip=2mu r'(d+n)$. When both $r$ and $r'$ scale with the model width $d$ (\emph{e.g.}, $\thickmuskip=2mu \medmuskip=2mu r=r'=\alpha d$ for a constant $\thickmuskip=2mu \medmuskip=2mu 0<\alpha\leq 1$), LoRA's parameter cost is $\mathcal{O}(d^2+dn)$, while OFT's is only $\mathcal{O}(d)$. To provide concrete numbers, consider a weight matrix of size $\thickmuskip=2mu \medmuskip=2mu 128\times 128$: LoRA introduces $2,048$ trainable parameters if $\thickmuskip=2mu \medmuskip=2mu r'=8$, whereas OFT requires just $960$ tunable parameters with $\thickmuskip=2mu \medmuskip=2mu r=8$ (assuming no block sharing).

\subsection{Constrained Orthogonal Finetuning}

OFT's expressivity can be further restrained by enforcing that the fine-tuned model remains close to the original, pretrained parameters. COFT realizes this by introducing a forward computation as follows:

\begin{equation}\label{eq:coft_general}
    \footnotesize
    \bm{z}=\bm{W}^\top\bm{x}=(\bm{R}\cdot\bm{W}^0)^\top\bm{x},~~~\text{s.t.}~\bm{R}^\top\bm{R}=\bm{R}\bm{R}^\top=\bm{I},~\norm{\bm{R}-\bm{I}}\leq \epsilon
\end{equation}

which imposes both an orthogonality condition and an $\epsilon$-ball constraint around the identity matrix. The orthogonality property may be realized via the Cayley parameterization as presented in Section~\ref{sect:eff_ortho}. However, ensuring the $\epsilon$-closeness constraint alongside the Cayley-based orthogonality is more intricate. To better understand the behavior of the Cayley transform, we consider the Neumann series to approximate $\thickmuskip=2mu \medmuskip=2mu \bm{R}=(\bm{I}+\bm{Q})(I-\bm{Q})^{-1}$ as $\thickmuskip=2mu \medmuskip=2mu \bm{R}\approx\bm{I}+2\bm{Q}+\mathcal{O}

Given that the series converges under the operator norm, it becomes possible to transfer the constraint $\thickmuskip=2mu \medmuskip=2mu\|\bm{R}-\bm{I}\|\leq\epsilon$ into the Cayley transform framework. This leads to a reformulated condition: $\thickmuskip=2mu \medmuskip=2mu\|\bm{Q}-\bm{0}\|\leq\epsilon'$, with $\bm{0}$ indicating a zero matrix and $\epsilon'$ representing a distinct error threshold from $\epsilon$. This revised condition imposed on $\bm{Q}$ can be straightforwardly satisfied via projected gradient descent techniques. To ensure the initial orthogonal matrix $\bm{R}$ starts as the identity, $\bm{Q}$ is simply initialized to the zero matrix. COFT effectively merges two explicit constraints: a limit on hyperspherical energy change, and a restriction on deviation from the pretrained parameters. While the latter is often an implicit regularization used in existing finetuning approaches, COFT treats it as a formalized constraint. Although OFT demonstrates strong results, our findings indicate that COFT enhances finetuning robustness further, attributable to the explicit control over parameter deviation. Figure~\ref{fig:eps_coft} showcases the influence of the $\epsilon$ parameter on COFT's effectiveness. From these results, it is clear that $\epsilon$ dictates the degree of adaptability during finetuning: larger $\epsilon$ values push the COFT outcomes closer to those of OFT, whereas reducing $\epsilon$ produces model behaviors that align more closely with the original pretrained text-to-image diffusion model.

\subsection{Re-scaled Orthogonal Finetuning}

To extend OFT, we introduce an additional trainable scaling factor for each neuron. This modification is based on the insight that changing the scale of a neuron does not affect its hyperspherical energy, since normalization will negate any changes in magnitude. In practice, the modified forward computation is: $\thickmuskip=2mu \medmuskip=2mu\bm{z}=(\bm{R}\bm{W}^0\bm{D})^\top\bm{x}$\footnote{\textbf{Errata}: The NeurIPS camera-ready version incorrectly reported the forward equation for re-scaled OFT as $\thickmuskip=2mu \medmuskip=2mu\bm{z}=(\bm{D}\bm{R}\bm{W}^0)^\top\bm{x}$. The original code, however, used the correct implementation, so the results in Appendix~\ref{app:rescaled} remain valid.} where $\thickmuskip=2mu \medmuskip=2mu \bm{D}=\text{diag}(s_1,\cdots,s_n)\in\mathbb{R}^{n\times n}$ is a parameterized diagonal matrix with each $s_1,\ldots,s_n$ constrained to be positive. Unlike the standard OFT as described in Eq.~\eqref{eq:oft_general}, which optimizes only over $\bm{R}$, this variation involves learning both $\bm{D}$ and the orthogonal matrix $\bm{R}$. The addition of $\bm{D}$ grants greater modeling capacity to OFT with minimal increase in parameter count. In our main experiments, we focus on the unscaled OFT to highlight the sole impact of orthogonal transformations, but our evaluations (see Appendix~\ref{app:rescaled}) indicate that incorporating the scaling layer generally yields improved results.

\section{Intriguing Insights and Discussions}

\textbf{OFT exhibits architecture-agnostic applicability}. The OFT approach is theoretically compatible with all neural network architectures. In the context of Transformers, LoRA is most commonly used for attention layer finetuning~\cite{hulora2022}, so for a direct performance comparison, we similarly restrict OFT to finetuning attention weights in our experiments. OFT is not only suitable for dense layers but is also naturally adapted for convolutional networks. Notably, the block-diagonal format of $\bm{R}$ offers interesting properties for convolution layers, distinguishing it from LoRA. For instance, if the number of blocks equals the number of input channels, each block manipulates an individual channel, functioning analogously to training separate depth-wise convolution filters~\cite{chollet2017xception}. Conversely, if all blocks are tied together, $\bm{R}$ applies an identical orthogonal transformation across all channels. Preliminary experiments applying OFT to convolutional layers are reported in Appendix~\ref{app:convolution}.

\textbf{Connection to LoRA}. LoRA introduces a low-rank matrix to ensure that the original pretrained weight matrix retains its core information by limiting dramatic changes. In contrast, OFT imposes an orthogonal (thus full-rank) transformation on the pretrained weight matrix, guaranteeing that the transformation preserves information learned during pretraining. The forward computation for OFT can be expressed as $\thickmuskip=2mu \medmuskip=2mu \bm{z}=(\bm{R}\bm{W}^0)^\top\bm{x}=(\bm{W}^0+(\bm{R}-\bm{I})\bm{W}^0)^\top\bm{x}$, where $\thickmuskip=2mu \medmuskip=2mu (\bm{R}-\bm{I})\bm{W}^0$ plays a similar role to the low-rank weight update used in LoRA. Typically, since $\bm{W}^0$ is full-rank, OFT effectively conducts a low-rank update as long as $\thickmuskip=2mu \medmuskip=2mu \bm{R}-\bm{I}$ has low rank. Like LoRA’s rank hyperparameter $r'$, OFT employs a block-diagonal parameter $r$ to constrain the count of trainable parameters. Importantly, LoRA and OFT exemplify two fundamentally different strategies for parameter efficiency: LoRA leverages low-rank decompositions, while OFT relies on a structured sparsity through block-diagonal orthogonal constraints.

\textbf{Why OFT converges faster}. Firstly, as shown in Figure~\ref{fig:toy_ae}, the most effective way to induce semantic change is by altering neuron orientations—precisely the mechanism of OFT. Moreover, OFT’s updates can be interpreted as performing fine-tuning on a smooth hyperspherical manifold, which offers more favorable optimization dynamics. This phenomenon is also supported by empirical findings in \cite{liu2021orthogonal}.

\textbf{Why not minimize hyperspherical energy}. Unlike the approach in \cite{liu2021orthogonal}, our objective is not to minimize hyperspherical energy. For classification tasks, it is important for neurons to be non-redundant, and while minimizing hyperspherical energy enforces a uniform distribution of neuron directions across the hypersphere, this regularization is not suitable for fine-tuning, as it risks discarding useful knowledge from pretraining.

\textbf{Balancing flexibility and regularity during finetuning}. Our investigation reveals an inherent tension between model flexibility and the need for regularity. Traditional finetuning techniques, though offering maximal flexibility, often compromise model stability and may even lead to model degeneration. Notably, OFT—a notably straightforward approach—strikes an effective balance by maintaining the pairwise angular relationships between neurons, which acts as an implicit regularization. The block-diagonal structure in parameterization can be interpreted as enforcing a more pronounced regularizing constraint on the orthogonal matrix.

\textbf{No increase in inference time}. In contrast to ControlNet, the OFT framework does not add any extra computational cost during inference. At inference, one can directly combine the learned orthogonal transformation $\bm{R}$ with the original pretrained weights $\bm{W}^0$, yielding a new, functionally equivalent parameter matrix $\thickmuskip=2mu \medmuskip=2mu \bm{W}=\bm{R}\bm{W}^0$. As a result, the runtime performance remains identical to that of the original model.

\section{Experiments and Results}

\textbf{Overall experimental protocol}. For our empirical evaluation, Stable Diffusion v1.5~\cite{rombach2022high} is selected as the backbone pretrained text-to-image generator. In order to maintain impartiality, samples from all methods are chosen randomly. For experiments focused on subject-driven generation, protocols follow those outlined by DreamBooth~\cite{ruiz2023dreambooth}, while experiments on controllable generation adopt procedures from ControlNet~\cite{zhang2023adding} and T2I-Adapter~\cite{mou2023t2i}. In making direct comparisons with LoRA, we restrict the application of OFT or COFT exclusively to the layers also targeted by LoRA. Additional details and experimental outcomes are provided in Appendix~\ref{app:settings}.

\subsection{Subject-driven Generation}

\textbf{Experimental details}. DreamBooth~\cite{ruiz2023dreambooth} and LoRA~\cite{hulora2022} serve as the primary baselines. The loss employed by all methods mirrors the one used in DreamBooth. For both DreamBooth and LoRA, the experimental procedure and hyperparameter choices adhere closely to their original publications. Supplementary results and more comprehensive information can be found in Appendix~\ref{app:settings},\ref{app:coft_oft},\ref{app:more_qualitative},\ref{app:failure}.

\textbf{Finetuning stability and convergence}. To assess the stability during finetuning and the rate at which convergence occurs, we compare DreamBooth, LoRA, OFT, and COFT. The relevant findings are displayed in Figure~\ref{fig:teaser} and Figure~\ref{fig:db_convergence}. It is evident that COFT and OFT consistently yield stable finetuning of the diffusion model. Within 400 iterations, DreamBooth and OFT methods effectively establish control over generation, whereas LoRA does not maintain subject identity. By 2000 iterations, DreamBooth begins producing degenerate outputs, and LoRA struggles to render the correct yellow shirt, instead generating yellow fur. Conversely, both OFT and COFT exhibit continued robustness and reliability in controlling image generation even after many iterations. These results provide evidence of OFT’s ability to combine rapid convergence with stable subject-driven generation. Importantly, OFT and COFT display a low sensitivity to the total number of finetuning steps, reducing the burden of hyperparameter selection across different subjects. This robustness allows users to choose a larger iteration count for OFT and COFT without laborious tuning. Moreover, when COFT employs a suitable $\epsilon$, both the learning rate and number of training iterations are straightforward to specify.

\textbf{Quantitative comparison}. In line with the protocol from \cite{ruiz2023dreambooth}, we carry out a quantitative analysis to measure subject fidelity (using DINO~\cite{caron2021emerging} and CLIP-I~\cite{radford2021learning}), prompt fidelity (CLIP-T~\cite{radford2021learning}), and the diversity of samples (LPIPS~\cite{zhang2018unreasonable}). For subject fidelity, CLIP-I reflects the mean pairwise cosine similarity between CLIP embeddings of real and generated images, while DINO follows a similar approach but relies on ViT S/16 DINO embeddings. Prompt fidelity, as measured by CLIP-T, represents the average cosine similarity between text prompts and their corresponding generated images. Sample diversity is quantified via the average LPIPS cosine similarity computed across generated outputs given identical subject and prompt. Results in Table~\ref{table:dreambooth_final} demonstrate that COFT and OFT notably surpass DreamBooth and LoRA on DINO and CLIP-I, while maintaining competitive or superior scores in prompt fidelity and diversity. To account for variability, each algorithm is finetuned on the same pretrained model with 30 separate random seeds. Overall, OFT is shown to not only improve convergence and stability, but also to deliver superior and consistent final results.

\textbf{Qualitative comparison}. To provide a clearer picture of the advantages offered by OFT, we present a selection of random samples from subject-driven generation tasks in Figure~\ref{fig:dreambooth_final}. In the interest of fairness, all samples are produced by applying each method to the same fine-tuned model, and no individualized text prompt tuning is performed for the final outputs. For every technique under evaluation, we utilize the model that achieves the highest validation CLIP scores. Examination of Figure~\ref{fig:dreambooth_final} reveals that both OFT and COFT consistently maintain strong semantic subject fidelity, whereas LoRA frequently does not preserve subject identity (for instance, in the bowl sample, LoRA entirely loses track of the subject). Additionally, OFT and COFT exhibit notably superior fidelity to the given text prompts. In contrast, although DreamBooth maintains subject identity, it is often unsuccessful in adhering to prompt instructions (see the first row for the bowl scenario). These findings from the qualitative analysis indicate that our OFT framework simultaneously attains enhanced controllability and subject integrity. Furthermore, OFT's performance is robust to the number of iterations—delivering stable results even as iterations increase, a property not shared by DreamBooth or LoRA. Supplementary qualitative cases are included in Appendix~\ref{app:more_qualitative}. A human subject study, detailed in Appendix~\ref{app:human}, provides additional evidence supporting the superiority of OFT.

\subsection{Controllable Generation}

\textbf{Settings}. For baseline comparisons, we employ ControlNet~\cite{zhang2023adding}, T2I-Adapter~\cite{mou2023t2i}, and LoRA~\cite{hulora2022}. The three challenging controllable generation tasks we focus on in this work are: Canny edge to image~(C2I) with the COCO dataset~\cite{lin2014microsoft}, segmentation map to image~(S2I) using the ADE20K dataset~\cite{zhou2017scene}, and landmark to face~(L2F) utilizing the CelebA-HQ dataset~\cite{karras2018progressive,xia2021tedigan}. Each approach is used to fine-tune Stable Diffusion~(SD) v1.5 for 20 epochs on these datasets. Additional experimental results are presented in Appendix~\ref{app:more_qualitative}, \ref{app:more_control_task}, and \ref{app:failure}.

\textbf{Convergence}. We investigate how quickly ControlNet, T2I-Adapter, LoRA, and COFT converge on the L2F task by conducting both quantitative and qualitative assessments. For our quantitative metric, we utilize the average $\ell_2$ distance between the control face landmarks and those predicted by the model. Figure~\ref{fig:face_convergence} presents the evolution of face landmark errors as a function of the number of fine-tuning epochs. Our observations indicate that both COFT and OFT display a markedly faster convergence rate. For example, LoRA requires 20 epochs to reach the same performance that OFT attains after just 8 epochs. It is worth emphasizing that OFT and COFT employ a number of trainable parameters comparable to LoRA, which remains substantially lower than the parameter count in ControlNet; despite this, they converge much more rapidly than existing techniques. Additional support for OFT's rapid convergence can be found in Figure~\ref{fig:teaser}, which illustrates that OFT is also notably more data-efficient than ControlNet and LoRA—achieving strong convergence using merely 5\% of the ADE20K dataset. Qualitatively, we primarily contrast OFT, COFT, and ControlNet, since ControlNet is closest to our approach in terms of landmark error. As depicted in Figure~\ref{fig:face_convergence_qual}, both OFT and COFT not only converge steadily but also progressively align the generated facial pose with the input control landmarks. Unlike our approaches, ControlNet exhibits a delayed emergence of controllability, which appears abruptly after the 8-th epoch; this abrupt transition coincides with the sudden convergence described by \cite{zhang2023adding}. The steady and smooth convergence exhibited by OFT enhances its practical utility, offering more reliable and predictable training dynamics. Consequently, tuning OFT’s hyperparameters becomes more straightforward.

\textbf{Quantitative comparison}. To assess the effectiveness of controllable generation, we propose a control consistency metric. This metric involves extracting the control signal from the output image and comparing it to the original input control signal. In the case of the C2I task, we utilize Intersection over Union (IoU) and F1 score for evaluation. For the S2I task, our assessment includes mean IoU as well as mean and overall accuracy. Regarding the L2F task, the evaluation is based on the average $\ell_2$ distance between the predicted landmarks and the control landmarks. Further elaboration on these consistency metrics can be found in Appendix~\ref{app:settings}. We select optimal hyperparameter configurations for each baseline in our comparison. The quantitative results in Table~\ref{table:control_gen} indicate that both OFT and COFT consistently achieve more robust and precise control than competing methods. Notably, approaches utilizing adapters (\emph{e.g.}, T2I-Adapter and ControlNet) tend to converge more slowly and generally produce inferior final outcomes. LoRA surpasses ControlNet in the S2I task, yet trails behind in the C2I and L2F tasks. Overall, our findings suggest that LoRA's maximum achievable performance remains limited, even with careful hyperparameter optimization. Conversely, the OFT framework continues to show improvements with increased numbers of trainable parameters, indicating it has not yet reached its performance limit. We highlight that the quantitative assessment methodology for controllable generation presented in this work represents a novel contribution, as it enables precise evaluation of the control fidelity of finetuned models, contingent on the accuracy of the underlying segmentation or detection models.

\textbf{Qualitative comparison}. We further provide a qualitative assessment of OFT and COFT in relation to leading approaches such as ControlNet, T2I-Adapter, and LoRA. As illustrated in the randomly sampled outputs in Figure~\ref{fig:control_final}, OFT and COFT consistently deliver images of high fidelity and realism, while maintaining precise controllability. Specifically, in the S2I task, LoRA fails to adhere to the input segmentation map, whereas ControlNet, OFT, and COFT successfully manage to guide the synthesis according to the segmentation. Notably, OFT and COFT surpass ControlNet by generating images with richer detail and more coherent geometric structures, all while utilizing fewer model parameters. For the C2I task, OFT and COFT demonstrate the ability to construct plausible and semantically meaningful images from sparse Canny edge inputs, outperforming both T2I-Adapter and LoRA. Regarding the L2F task, our method exhibits superior accuracy in pose control for synthesized faces, even in the case of extreme head poses. Across all three control scenarios, both OFT and COFT generate higher quality images than competing methods, substantiating the efficacy of the OFT framework in tasks requiring controlled generation. To further support this qualitative evaluation, we include additional examples for each control task in Appendix~\ref{app:control_exp}, and, in Appendix~\ref{app:more_control_task}, demonstrate the versatility of OFT on further control tasks such as dense pose to human body, sketch to image, and depth to image conversions.

\section{Concluding Remarks and Open Problems}

Inspired by the crucial role angular relationships among neurons play in shaping visual semantics, we introduce a straightforward yet powerful finetuning technique—orthogonal finetuning—for controlling text-to-image diffusion models. We apply this method to two key scenarios: subject-driven generation and controllable generation. In comparison to previous approaches, OFT offers enhanced control and more stable finetuning while utilizing fewer tunable parameters. Notably, OFT maintains computational efficiency during inference, thereby yielding models that are readily deployable.

Despite its advantages, OFT raises several noteworthy open questions. Firstly, while OFT achieves orthogonality through the Cayley parametrization, which relies on computing a matrix inverse, this can somewhat restrict its scalability. Although our block diagonal parametrization mitigates this issue, finding more efficient, differentiable approaches to this matrix inversion remains unresolved. Secondly, OFT enables an interesting form of compositionality, as orthogonal matrices generated from multiple OFT finetuning processes can be multiplied to yield another orthogonal matrix. Investigating whether these composite matrices can retain the specialized knowledge from each finetuning task is an appealing avenue for future research. Lastly, the parameter efficiency achieved by OFT heavily depends on the block diagonal structure; this construction, however, may introduce certain biases and constrain flexibility. Developing methods to enhance parameter efficiency in a manner that is both more effective and less biased constitutes another key challenge ahead.